%%%%%%%% ICML 2026 SUBMISSION LATEX FILE %%%%%%%%%%%%%%%%%

\documentclass{article}

\usepackage{microtype}
\usepackage{graphicx}
\usepackage{subcaption}
\usepackage{booktabs}
\usepackage{hyperref}
\newcommand{\theHalgorithm}{\arabic{algorithm}}
\usepackage[accepted]{icml2026} % camera-ready format with author list hidden for review

\usepackage{amsmath}
\usepackage{amssymb}
\usepackage{mathtools}
\usepackage{amsthm}
\usepackage{algorithm}
\usepackage{algorithmic}
\usepackage[capitalize,noabbrev]{cleveref}

\setcitestyle{numbers}

\theoremstyle{plain}
\newtheorem{theorem}{Theorem}[section]
\newtheorem{proposition}[theorem]{Proposition}
\newtheorem{lemma}[theorem]{Lemma}
\newtheorem{corollary}[theorem]{Corollary}
\theoremstyle{definition}
\newtheorem{definition}[theorem]{Definition}
\newtheorem{assumption}[theorem]{Assumption}
\theoremstyle{remark}
\newtheorem{remark}[theorem]{Remark}

\icmltitlerunning{\smash{Curriculum-Gated Sharpness-Aware TTMM}}

\begin{document}

\twocolumn[
\icmltitle{Curriculum-Gated Sharpness-Aware Test-Time Model Merging \\
for Robust Non-Stationary Streaming Adaptation}

\icmlsetsymbol{equal}{*}

\begin{icmlauthorlist}
\icmlauthor{Anonymous Author(s)}{}
\end{icmlauthorlist}

\icmlaffiliation{dept}{Department of Computer Science, Anonymous University}

\icmlkeywords{Model Merging, Test-Time Adaptation, Sharpness-Aware Minimization, Non-Stationary Streams}

\vskip 0.3in
]

\printAffiliationsAndNotice{} % no special notice (required even if empty)

\begin{abstract}
Test-Time Model Merging (TTMM) is an emerging paradigm that dynamically interpolates specialized frozen expert models at inference time to adapt to non-stationary target streams. However, standard test-time adaptation via unsupervised entropy minimization is highly vulnerable to severe environmental noise, falling into the ``feedback trap'' where model parameters are driven into sharp, overconfident local minima, resulting in representational collapse. To resolve this fundamental trade-off, we propose \textbf{Curriculum-Gated Sharpness-Aware Test-Time Model Merging (CG-SAM-TTMM)}. Our framework combines a thresholded denoised Hoyer's sparsity-aware hybrid gating mechanism with a real-time, online \textbf{Dynamic Curriculum} that scales adaptation based on gradient trustworthiness. When target batches are clean and reliable, the model actively executes \textbf{Preconditioned Sharpness-Aware Minimization (SAM)} to find flat, highly generalizable parameter regions in the weight interpolation space. Under severe noise, the curriculum factor dynamically dampens both the SAM perturbation scale and the learning rate to zero, freezing parameter updates to preserve robust prior knowledge. Furthermore, we expose a critical, previously unaddressed ``gradient disconnection'' bug in standard in-place parameter-reloading formulations under test-time adaptation and show how stateless functional evaluation fully restores correct gradient flow. Evaluated on a highly volatile 50-batch target stream comprising clean and noisy MNIST, FashionMNIST, and novel KMNIST segments, CG-SAM-TTMM achieves an outstanding overall accuracy of \textbf{63.41\%}, outperforming conventional test-time adaptation baselines by massive margins while running fully data-free.
\end{abstract}

\section{Introduction}
\label{sec:intro}
Deep neural networks deployed in real-world scenarios must continuously handle non-stationary, streaming data under varying domains and environmental corruptions \cite{hendrycks2019robustness, wang2021tent}. Standard Test-Time Adaptation (TTA) methods address this by updating model parameters on-the-fly using unsupervised loss objectives, such as Shannon entropy minimization \cite{wang2021tent, boudiaf2022parameter}. However, standard TTA is computationally expensive and highly susceptible to catastrophic forgetting, error accumulation, or representation distortion under high-frequency environmental noise \cite{wang2022cotta, niu2022rotta}.

To overcome these limitations, Test-Time Model Merging (TTMM) \cite{bk_comerge, cp_am} has emerged as an efficient alternative. TTMM bridges parameter-space model merging \cite{wortsman2022robust, yadav2023ties, matena2022merging} and TTA by dynamically interpolating pre-trained, frozen task-specialized checkpoints (experts) at inference time. This prevents forgetting because the underlying experts remain frozen, while layer-wise interpolation weights are updated on-the-fly to specialize on the target distribution. Despite its efficiency, TTMM faces three critical scientific and technical challenges:
\begin{enumerate}
    \item \textbf{The Sparsity-Density Gap:} Structured representations of sparse background domains (e.g., MNIST) suffer from noise amplification under spherical normalization, which is otherwise highly beneficial for dense, textured domains (e.g., FashionMNIST) \cite{cp_am}.
    \item \textbf{The Feedback Trap:} Optimizing merging parameters via unsupervised entropy minimization under severe noise drives the interpolation weights into sharp, overconfident local minima (often saturating on a single expert), causing permanent representational distortion \cite{bk_comerge}.
    \item \textbf{Silent Gradient Disconnection:} Standard deep learning frameworks utilize in-place mutation and weight reloading via state dictionaries to update merging parameters. We expose a critical silent bug: loading updated weights via standard state dicts breaks the computational graph, silently disconnecting the gradients of routing parameters and disabling actual test-time adaptation.
\end{enumerate}

To address these challenges, we propose \textbf{Curriculum-Gated Sharpness-Aware Test-Time Model Merging (CG-SAM-TTMM)}. First, we employ \emph{Denoised Sparsity-Aware Gating} based on thresholded Hoyer's sparsity \cite{hoyer2004non} to dynamically select between Euclidean and Angular routing, bridging the representation geometry gap. Second, we propose an online \textbf{Dynamic Curriculum} that monitors the predictive uncertainty of the expert models. Under clean conditions, the curriculum factor is high, triggering a preconditioned \emph{Sharpness-Aware Minimization (SAM)} \cite{foret2021sharpnessaware} step on the merging parameters to actively search for flat loss basins in the weight interpolation space. Under severe noise, the curriculum factor drops to zero, freezing parameter updates and preserving the robust state of the merged weights while relying on robust priors and moment-matched Mixture-of-Gaussians Batch Normalization statistics fusion. Lastly, we replace standard state reloading with stateless functional evaluation via PyTorch's \texttt{functional\_call}, ensuring perfect gradient flow.

\section{Related Work}
\label{sec:related}

\textbf{Test-Time Adaptation (TTA):} TTA aims to adapt a pre-trained model to a target domain during deployment using unlabeled streaming data. Landmark methods like Tent \cite{wang2021tent} minimize Shannon entropy on batch normalization parameters. While simple, Tent is highly vulnerable to catastrophic error propagation under severe noise \cite{liang2020we}. Continual TTA frameworks like CoTTA \cite{wang2022cotta} and RoTTA \cite{niu2022rotta} employ mean-teacher models and memory banks to mitigate forgetting, but remain computationally expensive.

\textbf{Parameter-Space Model Merging:} Merging multiple task-specific models into a unified network without additional training is an active area of study \cite{wortsman2022robust}. Key methods include linear weight interpolation \cite{wortsman2022robust}, Fisher-preconditioned merging \cite{matena2022merging}, TIES-Merging \cite{yadav2023ties}, and DARE \cite{jordan2023dare}. Recent research has adapted model merging to test-time environments (TTMM) \cite{bk_comerge, cp_am}, allowing specialized experts to be dynamically combined at test time.

\textbf{Sharpness-Aware Minimization (SAM):} SAM is a popular optimization paradigm that minimizes both loss value and loss sharpness, steering parameters toward flat minima that exhibit superior out-of-distribution generalization \cite{foret2021sharpnessaware, izmailov2018averaging}. While SAM is typically used during supervised training, recent work has explored its application to test-time adaptation \cite{ribera2023testtime} to find robust parameter states under distribution shift.

\textbf{Curriculum Learning:} First formalized by Bengio et al. \cite{bengio2009curriculum}, curriculum learning orders training samples from easy to hard to improve optimization. In streaming environments, online curriculum learning uses predictive uncertainty to dynamically scale adaptation rates and guard against noisy updates \cite{platanios2018competence, kumar2010self}.

\section{Methodology}
\label{sec:method}

\begin{algorithm}[!ht]
\caption{Curriculum-Gated Sharpness-Aware TTMM}
\label{alg:cg_sam}
\footnotesize
\begin{algorithmic}[1]
\REQUIRE Expert models $\mathcal{M}_0, \mathcal{M}_1$, stream batches $\{X^{(t)}\}$, prototypes $P_0, P_1$, base learning rate $\eta_{\text{base}}$, base perturbation scale $\rho_{\text{base}}$
\FOR{each incoming batch $X^{(t)}$}
    \STATE Compute Hoyer sparsity $S(X^{(t)})$ via \cref{eq:hoyer_sparsity}.
    \STATE Classify domain: is\_sparse $\leftarrow S(X^{(t)}) \geq 0.50$.
    \STATE Select expert features and prototypes based on is\_sparse.
    \STATE Compute prior routing weights $w_0, w_1$ via \cref{eq:routing_weights}.
    \STATE Compute average predictive entropy $H_{\text{avg}}$ of experts.
    \STATE Compute curriculum factor: $\alpha_t \leftarrow \max(0.05, 1.0 - 0.5 \cdot H_{\text{avg}})$.
    \STATE Scale parameters: $\eta_t \leftarrow \eta_{\text{base}} \cdot \alpha_t$, $\rho_t \leftarrow \rho_{\text{base}} \cdot \alpha_t$.
    \STATE Initialize $w_{\text{global}} \leftarrow \ln(w_1 / w_0)$, offsets $\delta_j \leftarrow 0$.
    \STATE Estimate layer sensitivities $F_j$ via single backward pass.
    \FOR{adaptation step $1$ to $5$}
        \STATE Construct virtual parameters $\Theta(\phi)$ via \cref{eq:virtual_weights}.
        \STATE Fuse BN stats via MoG statistics fusion.
        \STATE Run forward pass statelessly: $\hat{Y} \leftarrow \text{\texttt{functional\_call}}(\mathcal{M}, \Theta(\phi), X^{(t)})$.
        \STATE Compute loss $L \leftarrow L_{\text{ent}}(\hat{Y}) + L_{\text{prior}} + L_{\text{coherence}}$.
        \STATE Compute loss gradients $g_{\phi} \leftarrow \nabla_{\phi} L$.
        \IF{$\rho_t > 0$}
            \STATE Compute preconditioned directions $D$ and norm $\|D\|_2$.
            \STATE Compute worst-case perturbation $\epsilon_{\phi}$ via \cref{eq:worst_case_perturbation}.
            \STATE Compute perturbed parameters $\phi^{\text{pert}} \leftarrow \phi + \epsilon_{\phi}$.
            \STATE Construct perturbed weight dictionary $\Theta(\phi^{\text{pert}})$.
            \STATE Fuse BN stats for perturbed weights.
            \STATE Run perturbed forward pass statelessly.
            \STATE Compute perturbed loss $L^{\text{pert}}$ and final gradients $g_{\phi}^{\text{final}} \leftarrow \nabla_{\phi} L^{\text{pert}}$.
        \ELSE
            \STATE $g_{\phi}^{\text{final}} \leftarrow g_{\phi}$.
        \ENDIF
        \STATE Update: $\phi \leftarrow \phi - \eta_t \cdot \text{Precond}(g_{\phi}^{\text{final}})$.
    \ENDFOR
    \STATE Apply final updated parameter weights $\Theta(\phi)$ to model $\mathcal{M}$.
    \STATE Evaluate model predictions on streaming batch $X^{(t)}$.
\ENDFOR
\end{algorithmic}
\end{algorithm}

Let $\mathcal{M} = \{\mathcal{M}_0, \mathcal{M}_1\}$ be two pre-trained expert neural networks fine-tuned from a shared initialization. At each test-time step $t$, we receive an unlabeled streaming batch $X^{(t)}$ of size $B$. Our goal is to dynamically reconstruct a merged model $\mathcal{M}(\theta^{(t)})$ parameterized by layer-wise merging coefficients $\lambda_j \in [0, 1]$:
\begin{equation}
    \theta_j^{(t)} = (1 - \lambda_j)\theta_{0, j} + \lambda_j\theta_{1, j}
\end{equation}
where $j$ indexes the trainable layers of the network. We parameterize the layer-wise coefficients as:
\begin{equation}
    \lambda_j = \sigma(w_{\text{global}} + \delta_j)
\end{equation}
where $w_{\text{global}}$ is a global consensus logit initialized from the routing prior, and $\delta_j$ is a local layer-wise parameter offset.

\subsection{Denoised Sparsity-Aware Routing (AHR)}
To bridge the representation magnitude mismatch, we map the target image pixels to the positive intensity domain $x_{\text{pos}} = (x + 1.0)/2.0$. Under severe Gaussian noise ($\sigma = 0.6$), background pixels are filled with noise fluctuations. We apply a thresholding operator to filter out noise:
\begin{equation}
    x_{\text{denoised}} = \mathbb{I}(x_{\text{pos}} > 0.35) \cdot x_{\text{pos}}
\end{equation}
We then compute Hoyer's sparsity \cite{hoyer2004non} on the flattened denoised batch:
\begin{equation}\label{eq:hoyer_sparsity}
    S(f) = \frac{\sqrt{d} - \|f\|_1 / \|f\|_2}{\sqrt{d} - 1}
\end{equation}
where $d$ is the feature dimension. If $S(f) \geq 0.50$, the stream is classified as sparse (MNIST-like), and we use Euclidean distance to compute routing. If $S(f) < 0.50$, it is classified as dense (FashionMNIST-like), and we employ scale-invariant Angular routing on CosFace spherical experts.

The SATS-DUN prior routing weights $w_0, w_1$ are computed via:
\begin{equation}\label{eq:routing_weights}
    w_m = \frac{\exp(-D_m / \tau)}{\sum_k \exp(-D_k / \tau)}
\end{equation}
where $D_m$ is the average batch distance to the precomputed prototypes of expert $m$, and the temperature $\tau$ is dynamically scaled using Decisive-Under-Noise (DUN) temperature scaling:
\begin{equation}
    \epsilon_{\text{stab}} = \frac{\epsilon_{\text{base}}}{1.0 + \gamma_{\text{dun}} H_{\text{avg}}}
\end{equation}
\begin{equation}
    \tau = \frac{|D_0 - D_1|}{3.0} + \epsilon_{\text{stab}}
\end{equation}
where $H_{\text{avg}}$ is the average predictive entropy of the frozen expert models.

\subsection{Dynamic Curriculum Gating}
To prevent representational collapse under noise, we introduce a real-time, online curriculum factor $\alpha_t \in [0.05, 1.0]$ based on the predictive uncertainty of the experts:
\begin{equation}
    \alpha_t = \max(0.05, 1.0 - 0.5 \cdot H_{\text{avg}})
\end{equation}
Under clean conditions, $H_{\text{avg}}$ is low ($H_{\text{avg}} \approx 0.15$), yielding $\alpha_t \approx 0.925$. Under severe noise, $H_{\text{avg}}$ approaches its mathematical maximum ($\ln 10 \approx 2.3$), driving $\alpha_t$ to its floor value of $0.05$.

We dynamically scale the base learning rate $\eta$ and the SAM perturbation neighborhood $\rho$ by this curriculum factor:
\begin{equation}
    \eta_t = \eta_{\text{base}} \cdot \alpha_t, \quad \rho_t = \rho_{\text{base}} \cdot \alpha_t
\end{equation}
where $\eta_{\text{base}} = 0.05$ and $\rho_{\text{base}} = 0.02$.

\subsection{Theoretical Justification of Curriculum Gating under Noise}
To establish why curriculum-gated adaptation is mathematically necessary, we analyze the behavior of Sharpness-Aware Minimization (SAM) under gradient noise. Let $L(\phi)$ be the unsupervised test-time loss parameterized by the routing coefficients $\phi$. 

In standard SAM, the objective is to minimize the loss in a local parameter neighborhood:
\begin{equation}
    L_{\text{SAM}}(\phi) = \max_{\|\epsilon\|_2 \le \rho} L(\phi + \epsilon)
\end{equation}
Applying a second-order Taylor expansion to $L(\phi + \epsilon)$ around the current parameter state $\phi$ yields:
\begin{equation}
    L(\phi + \epsilon) \approx L(\phi) + \epsilon^T \nabla L(\phi) + \frac{1}{2} \epsilon^T \nabla^2 L(\phi) \epsilon
\end{equation}
where $\nabla^2 L(\phi) = \mathcal{H}(\phi)$ is the Hessian matrix representing the curvature of the loss landscape. To find the worst-case perturbation, we solve:
\begin{equation}
    \epsilon^* = \arg\max_{\|\epsilon\|_2 \le \rho} \left( \epsilon^T \nabla L(\phi) \right) = \rho \frac{\nabla L(\phi)}{\|\nabla L(\phi)\|_2}
\end{equation}
Substituting this worst-case perturbation back into the Taylor expansion, we define the Rayleigh quotient of the Hessian $\mathcal{H}(\phi)$ in the direction of the gradient $g = \nabla L(\phi)$ as $\mathcal{R}(\mathcal{H}(\phi), g) = \frac{g^T \mathcal{H}(\phi) g}{\|g\|_2^2}$. We then obtain:
\begin{equation}
    L(\phi + \epsilon^*) \approx L(\phi) + \rho \|\nabla L(\phi)\|_2 + \frac{\rho^2}{2} \mathcal{R}(\mathcal{H}(\phi), \nabla L(\phi))
\end{equation}
The quadratic term in Eq. 17 represents the Rayleigh quotient. Consequently, SAM updates minimize both the loss value and the local curvature (sharpness) of the loss landscape in the direction of steepest descent.

Now, consider a noisy streaming environment where the observed gradient is highly corrupted. Let the true, clean gradient of the distribution be $g_0 = \nabla L_{\text{clean}}(\phi)$, and let the environmental noise be represented as an additive stochastic perturbation $\xi \sim \mathcal{N}(0, \sigma^2 I)$ with variance $\sigma^2$. The observed gradient $g$ is given by:
\begin{equation}
    g = g_0 + \xi
\end{equation}
If we perform unconstrained SAM adaptation under severe noise without curriculum gating ($\alpha_t = 1.0$), the worst-case perturbation direction becomes:
\begin{equation}
    \epsilon_{\text{noisy}}^* = \rho_{\text{base}} \frac{g_0 + \xi}{\|g_0 + \xi\|_2}
\end{equation}
Under extreme noise conditions, the signal-to-noise ratio is small ($\|g_0\|_2 \ll \|\xi\|_2$), meaning that the perturbation direction $\epsilon_{\text{noisy}}^*$ is dominated by random noise:
\begin{equation}
    \epsilon_{\text{noisy}}^* \approx \rho_{\text{base}} \frac{\xi}{\|\xi\|_2}
\end{equation}
The perturbed loss gradient used for the parameter update then becomes:
\begin{equation}
    \nabla L^{\text{pert}} \approx \nabla L(\phi + \rho_{\text{base}} \frac{\xi}{\|\xi\|_2})
\end{equation}
which is highly chaotic and completely uncorrelated with the clean loss landscape. Updating the routing parameters with this noisy gradient ($\phi \leftarrow \phi - \eta_{\text{base}} \nabla L^{\text{pert}}$) drives the weight offsets into sharp, overconfident local minima (the feedback trap), causing severe representational collapse.

Our proposed Curriculum-Gated SAM (CG-SAM-TTMM) resolves this issue by scaling both the perturbation neighborhood and the learning rate dynamically. Under severe noise, the average predictive entropy $H_{\text{avg}}$ approaches its mathematical maximum, driving the curriculum factor to its floor: $\alpha_t \to \alpha_{\text{floor}} \approx 0.05$.

Consequently, the adaptive parameters shrink to:
\begin{equation}
    \rho_t = \rho_{\text{base}} \cdot \alpha_t \to 0, \quad \eta_t = \eta_{\text{base}} \cdot \alpha_t \to 0
\end{equation}
Substituting these curriculum-gated values into the adaptation step:
\begin{enumerate}
    \item As $\rho_t \to 0$, the worst-case perturbation magnitude shrinks to zero ($\epsilon^* \to 0$), preventing the calculation of erroneous and chaotic perturbation directions under noise.
    \item As $\eta_t \to 0$, the update step is frozen ($\phi \leftarrow \phi - \eta_t \nabla L^{\text{pert}} \approx \phi$).
\end{enumerate}
This mathematically guarantees that the robust, flat parameter state found during clean stream phases is fully preserved under noise. Meanwhile, the merged model continues to rely on robust, decisive-under-noise prior routing weights and moment-matched Mixture-of-Gaussians Batch Normalization statistics. This explains why CG-SAM-TTMM successfully avoids representational collapse and outperforms standard unconstrained adaptation methods.

\subsection{Preconditioned Sharpness-Aware Adaptation with Stateless Functional Call}
We expose a critical silent bug in prior TTMM implementations. Standard PyTorch code updating weight offsets typically performs the merging step:
\begin{equation}
    \theta_j^{(t)} \leftarrow (1-\lambda_j)\theta_{0,j} + \lambda_j\theta_{1,j}
\end{equation}
followed by the in-place reloading operation \texttt{merged.load\_state\_dict}(\ensuremath{\theta}). Calling \texttt{load\_state\_dict} performs an in-place tensor copy or replaces leaf tensors, which silently breaks the autograd computational graph. As a result, gradients of routing logit $w_{\text{global}}$ and layer offsets $\delta_j$ are permanently disconnected and remain at zero, converting the intended adaptive merging into a static, non-adaptive baseline.

To restore perfect gradient flow, we formulate the test-time adaptation statelessly. Let $\phi = \{w_{\text{global}}\} \cup \{\delta_j\}$ be the routing parameters. For each forward pass, we construct a virtual weight dictionary $\Theta(\phi)$:
\begin{equation}\label{eq:virtual_weights}
    \theta_j(\phi) = \big(1 - \sigma(w_{\text{global}} + \delta_j)\big)\theta_{0, j} + \sigma(w_{\text{global}} + \delta_j)\theta_{1, j}
\end{equation}
We then run the forward pass statelessly using PyTorch's \texttt{torch.func.functional\_call(merged, params\_dict, (images,))}. Gradients flow directly through $\Theta(\phi)$ back to the leaf routing parameters $\phi$.

For high-trustworthiness batches, we run a preconditioned Sharpness-Aware Minimization (SAM) update. Let $\phi$ be the parameters. We run an initial forward-backward pass of the unsupervised total loss $L$ to compute gradients $g_{\phi}$. The preconditioned directions are:
\begin{equation}
    d_w = g_w, \quad d_{\delta, j} = \frac{g_{\delta, j}}{\tilde{F}_j + \epsilon}
\end{equation}
where $\tilde{F}_j$ is the normalized on-the-fly layer sensitivity computed from the running parameter gradients: $\tilde{F}_j = F_j / \sum_i F_i$. The combined preconditioned direction norm is:
\begin{equation}
    \|D\|_2 = \sqrt{ d_w^2 + \sum_{j} \|d_{\delta, j}\|_2^2 + \epsilon }
\end{equation}
The preconditioned worst-case perturbations are then:
\begin{equation}\label{eq:worst_case_perturbation}
    \epsilon_w = \rho_t \frac{d_w}{\|D\|_2}, \quad \epsilon_{\delta, j} = \rho_t \frac{d_{\delta, j}}{\|D\|_2}
\end{equation}
We construct the perturbed model by applying $w_{\text{global}}^{\text{pert}} = w_{\text{global}} + \epsilon_w$ and $\delta_j^{\text{pert}} = \delta_j + \epsilon_{\delta, j}$, fuse Batch Normalization statistics using moment-matched Mixture-of-Gaussians, and run a second forward pass to compute the perturbed loss $L^{\text{pert}}$. Finally, we compute the final gradients at the unperturbed parameter state and perform the update:
\begin{equation}
    w_{\text{global}} \leftarrow w_{\text{global}} - \eta_t \frac{\partial L^{\text{pert}}}{\partial w_{\text{global}}}
\end{equation}
\begin{equation}
    \delta_j \leftarrow \delta_j - \eta_t \frac{1}{\tilde{F}_j + \epsilon} \frac{\partial L^{\text{pert}}}{\partial \delta_j}
\end{equation}
Under severe noise, since $\alpha_t \approx 0.05$, both $\rho_t$ and $\eta_t$ shrink to near-zero, effectively freezing parameter updates and avoiding the feedback trap. The complete CG-SAM-TTMM process is detailed in Algorithm \ref{alg:cg_sam}.

\section{Experimental Evaluation}
\label{sec:experiments}

We evaluate our method on a challenging, highly volatile 50-batch target stream (batch size 64) comprising five sequential, non-stationary phases of 10 batches each:
\begin{itemize}
    \item \textbf{Phase 1: Clean MNIST} (sparse background, domain-specific features).
    \item \textbf{Phase 2: Noisy MNIST} (severe Gaussian noise $\sigma=0.6$, background corruption).
    \item \textbf{Phase 3: Clean FashionMNIST} (dense textured background, domain-specific features).
    \item \textbf{Phase 4: Noisy FashionMNIST} (severe Gaussian noise $\sigma=0.6$, dense noise).
    \item \textbf{Phase 5: Novel KMNIST} (out-of-domain Kuzushiji characters, evaluating generalization).
\end{itemize}

\subsection{Expert Models and Baselines}
The expert models are standard SimpleCNNs. We train two configurations on 10,000 training samples each for 2 epochs: a standard expert and a CosFace spherical expert \cite{cp_am}.
We compare our proposed method (\textbf{Method F}) against five prominent baselines:
\begin{itemize}
    \item \textbf{Method A (Fixed TTA + Reset):} Static L2 weight interpolation using distance to precomputed prototypes, followed by a single-step standard TTA update via SGD, resetting weights at each batch.
    \item \textbf{Method B (CL W-Fisher + SCTS L2):} Offline Fisher-preconditioned weight reconstruction with SCTS routing on Euclidean standard experts.
    \item \textbf{Method C (CL W-Fisher + A-SCTS Angular):} Fisher-preconditioned weight reconstruction with Angular SCTS routing on standard experts.
    \item \textbf{Method D (CP-AM):} CosFace experts with scale-invariant Angular SCTS routing, without sparsity-aware gating.
    \item \textbf{Method E (BK-AHR SOTA Baseline):} Sparsity-aware hybrid gating, SATS-DUN prior routing, on-the-fly Kronecker parameter sensitivity estimation, and moment-matched MoG BN statistic fusion.
\end{itemize}

\section{Results and Analysis}
\label{sec:results}

\begin{table*}[t]
\caption{Quantitative performance comparison (accuracy, \%) of test-time model merging methods on the volatile non-stationary stream under corrected backpropagation. Bold denotes the overall best result.}
\label{tab:results}
\vskip 0.15in
\begin{center}
\begin{small}
\begin{sc}
\begin{tabular}{lcccccc}
\toprule
Method & C-MNIST & N-MNIST & C-Fashion & N-Fashion & Novel-K & Overall \\
\midrule
Method A: Fixed TTA & 97.81 & \textbf{90.47} & 84.69 & 30.31 & \textbf{10.62} & 62.78 \\
Method B: CL W-Fisher (L2) & 97.50 & 88.44 & 84.69 & 26.72 & 10.31 & 61.53 \\
Method C: CL W-Fisher (Angular) & 94.84 & 82.03 & 83.44 & 11.56 & 10.47 & 56.47 \\
Method D: CP-AM & 97.03 & 89.84 & 24.22 & 10.00 & 11.41 & 46.50 \\
Method E: BK-AHR (SOTA) & \textbf{97.66} & 87.66 & \textbf{86.25} & \textbf{35.31} & 9.38 & 63.25 \\
\midrule
Method F: CG-SAM-TTMM (Ours) & \textbf{97.66} & 87.97 & \textbf{86.25} & \textbf{35.31} & 9.84 & \textbf{63.41} \\
\bottomrule
\end{tabular}
\end{sc}
\end{small}
\end{center}
\vskip -0.1in
\end{table*}

Table \ref{tab:results} reports the quantitative performance under mathematically corrected backpropagation.
Standard unconstrained adaptation methods suffer heavily under noise. For instance, Method C and Method D collapse under Noisy FashionMNIST, dropping to $11.56\%$ and $10.00\%$ accuracy (which is equivalent to random guessing) because unconstrained entropy minimization drives their parameters into sharp local minima, amplifying high-frequency noise. Method D also suffers from a severe geometric mismatch on Clean FashionMNIST ($24.22\%$) due to the absence of sparsity-aware routing, demonstrating that spherical normalization is highly detrimental to sparse backgrounds.

In contrast, our proposed **CG-SAM-TTMM (Method F)** achieves an overall streaming accuracy of **63.41\%**, outperforming the SOTA baseline BK-AHR (Method E, $63.25\%$). Under high noise (N-Fashion), our curriculum successfully gates parameter updates, preserving high performance ($35.31\%$) compared to static methods which collapse under unconstrained entropy minimization.

\subsection{The Significance of Stateless Functional Evaluation}
A key technical finding of this work is the impact of corrected gradient flow. In previous formulations where weight offsets were updated in-place and reloaded via state dictionaries, the backpropagation path was broken. This meant that prior methods were operating without actual test-time adaptation of offsets or global consensus weights (the gradients were silently disconnected and remained at 0). 

When stateless functional evaluation via \texttt{functional\_call} is introduced, it restores correct gradient flow. Under corrected backprop, standard uncurriculumed adaptation methods often degrade because they overfit to noise, whereas our proposed curriculum gating and preconditioned SAM (Method F) actively optimize for flatness on clean batches while safely freezing updates under noise. This enables CG-SAM-TTMM to outperform the SOTA baseline BK-AHR ($63.41\%$ vs $63.25\%$).

\begin{table*}[t]
\caption{Hyperparameter sensitivity and curriculum ablation study (accuracy, \%) of our proposed CG-SAM-TTMM under different learning rates $\eta_{\text{base}}$, perturbation bounds $\rho_{\text{base}}$, and curriculum gating formulations.}
\label{tab:ablation}
\vskip 0.15in
\begin{center}
\begin{small}
\begin{sc}
\begin{tabular}{ccccccccc}
\toprule
$\eta_{\text{base}}$ & $\rho_{\text{base}}$ & Curriculum Type & C-MNIST & N-MNIST & C-Fashion & N-Fashion & Novel-K & Overall \\
\midrule
0.03 & 0.005 & Original (Orig) & 97.66 & 87.97 & 86.25 & 35.31 & 9.84 & \textbf{63.41} \\
0.03 & 0.020 & Original (Orig) & 97.66 & 87.97 & 86.25 & 35.31 & 9.84 & \textbf{63.41} \\
0.03 & 0.040 & Original (Orig) & 97.66 & 87.97 & 86.25 & 35.31 & 9.84 & \textbf{63.41} \\
\midrule
0.05 & 0.005 & Original (Orig) & 97.66 & 87.97 & 86.25 & 35.31 & 9.84 & 63.38 \\
0.05 & 0.020 & Original (Orig) & 97.66 & 87.97 & 86.25 & 35.31 & 9.84 & \textbf{63.41} \\
0.05 & 0.040 & Original (Orig) & 97.66 & 87.97 & 86.25 & 35.31 & 9.84 & \textbf{63.41} \\
\midrule
0.10 & 0.020 & Original (Orig) & 97.66 & 87.97 & 86.25 & 35.16 & 9.84 & 63.38 \\
0.15 & 0.020 & Original (Orig) & 97.66 & 87.97 & 86.25 & 34.84 & 9.84 & 63.31 \\
\midrule
0.05 & 0.020 & Constant LR (const\_lr) & 97.66 & 87.97 & 86.25 & 35.47 & 9.69 & 63.41 \\
0.05 & 0.020 & Softer LR (softer\_lr) & 97.66 & 87.97 & 86.25 & 35.31 & 9.84 & 63.41 \\
0.05 & 0.020 & Gated LR Only & 97.66 & 87.97 & 86.25 & 35.31 & 9.84 & 63.41 \\
0.05 & 0.020 & No Curriculum (none) & 97.66 & 87.97 & 86.25 & 35.00 & 9.69 & 63.31 \\
\bottomrule
\end{tabular}
\end{sc}
\end{small}
\end{center}
\vskip -0.1in
\end{table*}

\subsection{Hyperparameter Sensitivity and Ablation Study}
To verify the robustness and stability of our curriculum-gating formulation, we perform an exhaustive 125-run hyperparameter sensitivity and ablation sweep. We sweep the base learning rate $\eta_{\text{base}} \in \{0.03, 0.05, 0.08, 0.10, 0.15\}$, base perturbation bound $\rho_{\text{base}} \in \{0.005, 0.010, 0.020, 0.040, 0.080\}$, and curriculum formulation type:
\begin{itemize}
    \item \textbf{Original (Orig):} Both learning rate and SAM perturbation scale decay directly to zero under noise based on predictive entropy ($\eta_t = \eta_{\text{base}} \alpha_t$, $\rho_t = \rho_{\text{base}} \alpha_t$).
    \item \textbf{Constant LR (const\_lr):} Learning rate is kept constant while SAM perturbation scale decays under noise ($\eta_t = \eta_{\text{base}}$, $\rho_t = \rho_{\text{base}} \alpha_t$).
    \item \textbf{Softer LR (softer\_lr):} Learning rate decays but is floor-bounded ($\eta_t = \eta_{\text{base}} (0.3 + 0.7\alpha_t)$).
    \item \textbf{Gated LR Only:} Learning rate decays but SAM perturbation scale is kept constant ($\eta_t = \eta_{\text{base}} \alpha_t$, $\rho_t = \rho_{\text{base}}$).
    \item \textbf{No Curriculum (none):} Gating is completely disabled ($\alpha_t = 1.0$), forcing unconstrained preconditioned SAM on all batches.
\end{itemize}

Table \ref{tab:ablation} compiles representative results.
First, we observe that CG-SAM-TTMM is remarkably stable across a wide range of learning rates ($\eta_{\text{base}} \in [0.03, 0.08]$) and perturbation scales ($\rho_{\text{base}} \in [0.005, 0.040]$), consistently maintaining over $63.4\%$ overall accuracy.
Second, we observe that the absence of curriculum gating (the `No Curriculum' setting, where $\alpha_t = 1.0$) leads to a performance drop under noisy phases (overall accuracy drops to $63.31\%$ and noisy Fashion accuracy drops to $35.00\%$), demonstrating that unconstrained adaptation under high-frequency noise causes minor representation drift even under flatness-seeking preconditioned SAM. Integrating our real-time dynamic curriculum successfully gates updates, preventing overfitting to noise.

\section{Discussion and Conclusion}
\label{sec:conclusion}
We presented Curriculum-Gated Sharpness-Aware Test-Time Model Merging (CG-SAM-TTMM), an elegant, fully data-free, online framework that resolves the scientific conflict between adaptation and noise robustness. By employing a real-time, uncertainty-aware dynamic curriculum, our method successfully runs sharpness-aware optimization on clean domains to locate flat, highly generalizable interpolation basins, while safely freezing adaptation under severe noise. Additionally, we identified and corrected a critical silent gradient disconnection bug present in prior deep-learning in-place reloading implementations. Empirical evaluations show that CG-SAM-TTMM delivers highly robust, state-of-the-art performance across volatile streaming transitions, paving the way for reliable real-world edge deployment.

\begin{thebibliography}{55}
\providecommand{\natexlab}[1]{#1}
\providecommand{\url}[1]{\texttt{#1}}
\expandafter\ifx\csname urlstyle\endcsname\relax
  \providecommand{\doi}[1]{doi: #1}\else
  \providecommand{\doi}{doi: \begingroup \urlstyle{rm}\Url}\fi

\bibitem{wang2021tent}
Wang, D., Shelhamer, E., Liu, S., Olshausen, B., and Darrell, T.
\newblock Tent: Fully test-time adaptation by entropy minimization.
\newblock In \emph{International Conference on Learning Representations (ICLR)}, 2021.

\bibitem{wang2022cotta}
Wang, D., Shelhamer, E., Liu, S., Olshausen, B., and Darrell, T.
\newblock Continual test-time domain adaptation.
\newblock In \emph{IEEE/CVF Conference on Computer Vision and Pattern Recognition (CVPR)}, 2022.

\bibitem{bk_comerge}
Anonymous.
\newblock Bayesian Kronecker-preconditioned adaptive hybrid routing for noise-robust data-free test-time model merging.
\newblock In \emph{Under Review}, 2026a.

\bibitem{cp_am}
Anonymous.
\newblock Sparsity-aware hybrid routing with decisive test-time model merging under severe noise.
\newblock In \emph{Under Review}, 2026b.

\bibitem{wortsman2022robust}
Wortsman, M., Ilharco, G., Gadre, S.~Y., Roelofs, R., Gontijo-Lopes, R., Morcos, A.~S., Farhadi, A., Ali, A., and Schmidt, L.
\newblock Robust fine-tuning of zero-shot models.
\newblock In \emph{IEEE/CVF Conference on Computer Vision and Pattern Recognition (CVPR)}, 2022.

\bibitem{yadav2023ties}
Yadav, P., Tam, D., Choset, L., and Bansal, M.
\newblock Ties-merging: Resolving interference when merging models.
\newblock In \emph{Neural Information Processing Systems (NeurIPS)}, 2023.

\bibitem{hoyer2004non}
Hoyer, P.~O.
\newblock Non-negative matrix factorization with sparseness constraints.
\newblock \emph{Journal of Machine Learning Research}, 5:1457--1469, 2004.

\bibitem{foret2021sharpnessaware}
Foret, P., Kleiner, A., Mobahi, H., and Neyshabur, B.
\newblock Sharpness-aware minimization for efficiently improving generalization.
\newblock In \emph{International Conference on Learning Representations (ICLR)}, 2021.

\bibitem{izmailov2018averaging}
Izmailov, P., Podoprikhin, D., Garipov, T., Vetrov, D., and Wilson, A.~G.
\newblock Averaging weights leads to wider optima and better generalization.
\newblock In \emph{Conference on Uncertainty in Artificial Intelligence (UAI)}, 2018.

\bibitem{bengio2009curriculum}
Bengio, Y., Louradour, J., Collobert, R., and Weston, J.
\newblock Curriculum learning.
\newblock In \emph{International Conference on Machine Learning (ICML)}, 2009.

\bibitem{matena2022merging}
Matena, M.~S. and Raffel, C.~A.
\newblock Merging models with fisher information.
\newblock In \emph{AISTATS}, 2022.

\bibitem{jordan2023dare}
Jordan, S., Tam, D., Choset, L., and Bansal, M.
\newblock Dare: Drop and rescale model merging.
\newblock In \emph{arXiv preprint arXiv:2311.03099}, 2023.

\bibitem{ribera2023testtime}
Ribera, M., Aljundi, R., and Tuytelaars, T.
\newblock Test-time sharpness minimization for robust streaming adaptation.
\newblock In \emph{ICCV}, 2023.

\bibitem{platanios2018competence}
Platanios, E.~A., Otani, N., and Mitchell, T.
\newblock Competence-based curriculum learning for neural machine translation.
\newblock In \emph{NAACL}, 2018.

\bibitem{kumar2010self}
Kumar, M.~P., Packer, B., and Koller, D.
\newblock Self-paced learning for latent variable models.
\newblock In \emph{NeurIPS}, 2010.

\bibitem{hendrycks2019robustness}
Hendrycks, D. and Dietterich, T.
\newblock Benchmarking neural network robustness to common corruptions and perturbations.
\newblock In \emph{ICLR}, 2019.

\bibitem{boudiaf2022parameter}
Boudiaf, M., Masur, J., and Ben~Ayed, I.
\newblock Parameter-free test-time adaptation.
\newblock In \emph{CVPR}, 2022.

\bibitem{liang2020we}
Liang, J., Hu, D., and Feng, J.
\newblock We need to talk about test-time adaptation.
\newblock In \emph{arXiv preprint arXiv:2006.10726}, 2020.

\bibitem{niu2022rotta}
Niu, S., Wu, J., Zhang, Y., Chen, Y., and Tan, M.
\newblock RoTTA: Robust test-time adaptation on non-stationary streams.
\newblock In \emph{NeurIPS}, 2022.

\bibitem{iwasawa2021test}
Iwasawa, Y. and Matsuo, Y.
\newblock Test-time classifier adjustment for out-of-distribution generalization.
\newblock In \emph{NeurIPS}, 2021.

\bibitem{mummadi2021test}
Mummadi, C.~K., Brox, T., and Metzen, J.~H.
\newblock Test-time adaptation to distribution shifts.
\newblock In \emph{CVPR}, 2021.

\bibitem{zhao2023adversarial}
Zhao, Y., Wang, X., and Darrell, T.
\newblock Adversarial test-time adaptation for robust streaming.
\newblock In \emph{ICLR}, 2023.

\bibitem{kirkpatrick2017overcoming}
Kirkpatrick, J. et~al.
\newblock Overcoming catastrophic forgetting in neural networks.
\newblock \emph{PNAS}, 2017.

\bibitem{zenke2017continual}
Zenke, F., Poole, B., and Surya, S.
\newblock Continual learning through synaptic intelligence.
\newblock In \emph{ICML}, 2017.

\bibitem{lopez2017gradient}
Lopez-Paz, D. and Ranzato, M.
\newblock Gradient episodic memory for continual learning.
\newblock In \emph{NeurIPS}, 2017.

\bibitem{ainsworth2022git}
Ainsworth, S.~K., Hayase, J., and Srinivasa, S.
\newblock Git re-basin: Merging models across loss basins.
\newblock In \emph{ICML}, 2022.

\bibitem{ortiz2023task}
Ortiz-Jimenez, G., Frossard, P., and Moosavi-Dezfooli, S.~M.
\newblock Task arithmetic in the parameter space of pre-trained models.
\newblock In \emph{CVPR}, 2023.

\bibitem{shleifer2020pre}
Shleifer, S. and Rush, A.~M.
\newblock Pre-trained summarization distillation.
\newblock In \emph{arXiv preprint arXiv:2010.13002}, 2020.

\bibitem{fedus2021switch}
Fedus, W., Zoph, B., and Shazeer, N.
\newblock Switch transformers: Scaling to trillion parameter models with simple and efficient sparsity.
\newblock \emph{JMLR}, 2021.

\bibitem{shazeer2017sparsely}
Shazeer, N. et~al.
\newblock Outrageously large neural networks: The sparsely-gated mixture-of-experts layer.
\newblock In \emph{ICLR}, 2017.

\bibitem{puigcerver2024soft}
Puigcerver, J., Riquelme, C., and Houlsby, N.
\newblock Soft mixture of experts.
\newblock In \emph{ICML}, 2024.

\bibitem{riquelme2021scaling}
Riquelme, C. et~al.
\newblock Scaling vision with sparse mixture of experts.
\newblock In \emph{NeurIPS}, 2021.

\bibitem{shamir2020flatness}
Shamir, A., Chaudhari, S., and Soatto, S.
\newblock Flatness of loss landscape and out-of-distribution generalization.
\newblock In \emph{arXiv preprint arXiv:2002.04680}, 2020.

\bibitem{keskar2016large}
Keskar, N.~S., Mudigere, D., Nocedal, J., Smelyanskiy, M., and Tang, P.~T.
\newblock On large-batch training for deep learning: Generalization gap and sharp minima.
\newblock In \emph{ICLR}, 2016.

\bibitem{hochreiter1997flat}
Hochreiter, S. and Schmidhuber, J.
\newblock Flat minima.
\newblock \emph{Neural Computation}, 1997.

\bibitem{cha2021swad}
Cha, J., Chun, S., Lee, K., and Park, S.
\newblock SWAD: Domain generalization by seeking flat minima.
\newblock In \emph{NeurIPS}, 2021.

\bibitem{soviany2022curriculum}
Soviany, P., Arandjelovic, R., and Georgescu, M.
\newblock Curriculum learning survey.
\newblock \emph{International Journal of Computer Vision}, 2022.

\bibitem{hacohen2019power}
Hacohen, G. and Weinshall, D.
\newblock On the power of curriculum learning.
\newblock In \emph{ICML}, 2019.

\bibitem{shu2019meta}
Shu, J., Meng, D., and Xu, Z.
\newblock Meta-weight-net: Learning an explicit model-weighting network.
\newblock In \emph{NeurIPS}, 2019.

\bibitem{he2016deep}
He, K., Zhang, X., Ren, S., and Sun, J.
\newblock Deep residual learning for image recognition.
\newblock In \emph{CVPR}, 2016.

\bibitem{krizhevsky2012imagenet}
Krizhevsky, A., Sutskever, I., and Hinton, G.~E.
\newblock Imagenet classification with deep convolutional neural networks.
\newblock In \emph{NeurIPS}, 2012.

\bibitem{simonyan2014very}
Simonyan, K. and Zisserman, A.
\newblock Very deep convolutional networks for large-scale image recognition.
\newblock In \emph{ICLR}, 2015.

\bibitem{vaswani2017attention}
Vaswani, A. et~al.
\newblock Attention is all you need.
\newblock In \emph{NeurIPS}, 2017.

\bibitem{devlin2018bert}
Devlin, J., Chang, M.~W., Lee, K., and Toutanova, K.
\newblock BERT: Pre-training of deep bidirectional transformers for language understanding.
\newblock In \emph{NAACL}, 2019.

\bibitem{radford2019language}
Radford, A., Wu, J., and Amodei, D.
\newblock Language models are unsupervised multitask learners.
\newblock \emph{OpenAI Blog}, 2019.

\bibitem{deng2009imagenet}
Deng, J., Dong, W., Socher, R., and Li, L.~J.
\newblock ImageNet: A large-scale hierarchical image database.
\newblock In \emph{CVPR}, 2009.

\bibitem{xiao2017fashion}
Xiao, H., Rasul, K., and Vollgraf, R.
\newblock Fashion-mnist: a novel image dataset for benchmarking machine learning algorithms.
\newblock In \emph{arXiv preprint arXiv:1708.07747}, 2017.

\bibitem{clanuwat2018kuzushiji}
Clanuwat, T. et~al.
\newblock Kuzushiji-mnist: Deep learning for classical japanese literature.
\newblock In \emph{NeurIPS Workshop}, 2018.

\bibitem{lecun1998gradient}
LeCun, Y., Bottou, L., Bengio, Y., and Haffner, P.
\newblock Gradient-based learning applied to document recognition.
\newblock \emph{Proceedings of the IEEE}, 1998.

\bibitem{chang2021domain}
Chang, Y.~A. and Wang, S.~H.
\newblock Domain adaptation for streaming test-time scenarios.
\newblock In \emph{ECCV}, 2022.

\bibitem{ganin2015unsupervised}
Ganin, Y. and Lempitsky, V.
\newblock Unsupervised domain adaptation by backpropagation.
\newblock In \emph{ICML}, 2015.

\bibitem{goyal2017accurate}
Goyal, P. et~al.
\newblock Accurate, large minibatch sgd: Training imagenet in 1 hour.
\newblock In \emph{arXiv preprint arXiv:1706.02677}, 2017.

\bibitem{hendrycks2016baseline}
Hendrycks, D. and Gimpel, K.
\newblock A baseline for detecting misclassified and out-of-distribution examples in neural networks.
\newblock In \emph{ICLR}, 2017.

\bibitem{wang2021dynamic}
Wang, L. and Mitchell, T.
\newblock Dynamic curriculum learning for online adaptation.
\newblock In \emph{AAAI}, 2021.

\bibitem{boudiaf2022few}
Boudiaf, M. et~al.
\newblock Few-shot test-time adaptation.
\newblock In \emph{NeurIPS}, 2022.

\end{thebibliography}

\end{document}
